\HMTNarrativeHeading{La forma conservada por la segunda proyección}

La realización excepcional desarrolla la incidencia que acompaña a
la generación numérica. Cada prefijo transmite relaciones de frontera
compatibles con los truncamientos. La codificación de Paley conserva
la palabra y su imagen; el diseño de Witt organiza sus soportes; la
elevación hacia Golay introduce la incidencia binaria con sus marcas.
El recorrido posterior construye las realizaciones reticulares y
operatoriales correspondientes.

Las estrellas de incidencia permiten observar este enlace. El
registro de acoplamiento determina una cuaterna. Las hexadas que la
contienen distinguen cuatro pares complementarios y una firma
selecciona la hexada marcada. En la figura, las uniones representan
operaciones entre conjuntos. La elevación a octadas conserva esa
configuración y distingue una quinta octada transversal. Las cinco
regiones numéricas mantienen así una correspondencia con relaciones
de incidencia especificadas.

La igualdad de una lectura ternaria puede reunir regiones cuya
orientación, multiplicidad y posición siguen siendo distintas. Esos
datos intervienen en la correspondencia excepcional. La segunda
proyección conserva precisamente relaciones que la lectura de valor
trata de otro modo. La compatibilidad entre ambas se expresa sobre
el estado enriquecido común y sus mapas de refinamiento.

La teoría M y la dualidad T se presentan mediante las realizaciones
dimensionales y sus condiciones de compatibilidad. La acción del
Monstruo tiene su portador en el espacio construido para ella. Los
lectores numéricos, las configuraciones finitas, los retículos y las
representaciones reciben dominios diferenciados dentro de una misma
genealogía. Esta precisión permite seguir cómo una estructura
incidencial se transmite hasta sus consecuencias posteriores.

La parte siguiente conserva las construcciones, tablas y pruebas
que hacen efectivo el recorrido. La narración devuelve cada
realización a la cuestión rectora: qué relaciones sobreviven a la
proyección y cómo se recuperan en el objeto al que llegan. La
conservación de una forma de Gram, la covarianza de una codificación
y la invariancia de una acción tendrán pruebas propias. Su
articulación aporta el contenido geométrico de la conservación
evolutiva de la unidad.

